\documentclass[10pt,a4paper]{amsart}
\usepackage[margin=19mm]{geometry}
\usepackage{amsmath,amssymb,mathtools}
\usepackage[scr=rsfs,cal=euler]{mathalfa}
\usepackage{xfrac}
\usepackage[unicode,hidelinks,bookmarks=false]{hyperref}
\newcommand{\ie}{i.e.\ }
\newcommand{\cf}{cf.\ }
\newcommand{\resp}{respectively\ }
\newcommand{\Subsection}{\subsection}
\providecommand{\nobreakdash}{\nobreak\hbox{-}}
\setlength{\emergencystretch}{2em}
\multlinegap0pt
\usepackage{bm}
\usepackage{bbm}
\usepackage{dsfont}
\usepackage{xspace}
\usepackage{cancel}
\usepackage{mathtools}
\usepackage{amsthm}
\makeatletter
\@ifundefined{theorem}{\newtheorem{theorem}{Theorem}}{}
\@ifundefined{definition}{\newtheorem{definition}[theorem]{Definition}}{}
\@ifundefined{lemma}{\newtheorem{lemma}[theorem]{Lemma}}{}
\@ifundefined{proposition}{\newtheorem{proposition}[theorem]{Proposition}}{}
\makeatother
\title[Sharp bounds for stochastic proximal and projection estimato]{Sharp bounds for stochastic proximal and projection estimators via radial dominance}
\author{Gonzalo Contador, Pedro P\'erez-Aros, Emilio Vilches}
\date{Source: arXiv:2607.08670v2 (CC BY 4.0)}
\begin{document}
\maketitle

\noindent\emph{Source and licence.} Sharp bounds for stochastic proximal and projection estimators via radial dominance. Version arXiv:2607.08670v2 (CC BY 4.0), distributed under the Creative Commons Attribution 4.0 International licence, as recorded in the arXiv licence metadata for this exact version and shown on the abstract page https://arxiv.org/abs/2607.08670v2. Full rights notice: licenses/arxiv-2607.08670-CC-BY-4.0.txt.\par\smallskip

\noindent\textit{Editorial scope.} Counted unit: the Proposition whose source label is \texttt{stochastic-dominance-lemma} in the article (source0.tex lines 628--639, source label \texttt{stochastic-dominance-lemma}), delivered together with the article's own complete original proof of it (source0.tex lines 640--689, one \texttt{proof} environment of 50 lines). No mathematics is written, repaired or re-derived: statement and proof are the article's own text line for line. Required context retained uncounted: Admissible (lines 489--507); dominance (lines 503--505); Stochastic-dominance (lines 574--601); polar-decomposition (lines 605--613); likelihood-ratio (lines 615--624). Every definition, display and label of those lines is retained unchanged. Omitted from this package and not redistributed: 1--488, 506--573, 602--604, 614--614, 625--627, 690--1669. Nothing inside a retained excerpt is omitted or altered: every retained body is delivered exactly as the article's lines are written, so no omission receipt of any kind accompanies this package. Citations: the retained excerpts carry no citation command at all, so no bibliography entry of the article is transcribed and no reference is redistributed. Reference adaptation: none was needed and none was made; every reference inside the delivered unit resolves inside it, so no unresolved reference or citation is produced. Printed slips kept literal: no printed word, symbol, hypothesis or number of the counted statement or of its proof was altered. Layout and class: the article's own class is \texttt{article} with packages including geometry, fontenc, inputenc, amssymb, mathtools, microtype, enumitem, subcaption, xcolor, booktabs, pgfplots, hyperref, dsfont, mathrsfs; the wrapper is the lane's pinned \texttt{amsart} editorial wrapper at 10pt on A4, which restates the theorem-like environments the retained text needs and adds the article's own macro definitions that the retained text needs (none), with extra wrapper packages (bm, bbm, dsfont, xspace, cancel, mathtools). The delivered numbering is the unit's own. Assets: no retained span contains a figure environment, a graphics-inclusion command, a PDF-page inclusion, a table or a TikZ picture; no image file is copied into this package, the package declares no asset dependency and no image file is redistributed with it. Licence: version arXiv:2607.08670v2 of this article is distributed under the Creative Commons Attribution 4.0 International licence, as recorded in the arXiv licence metadata for this exact version and shown on the abstract page https://arxiv.org/abs/2607.08670v2. A full rights notice is delivered at licenses/arxiv-2607.08670-CC-BY-4.0.txt. Characters: every retained excerpt is pure ASCII, so the delivered engine, XeLaTeX with its default Latin Modern text font, reports no missing character.
\begin{definition}Let $G\colon \mathbb{R}^n \to \mathbb{R}\cup \{+\infty\}$ be proper and lower semicontinuous, with $G(0)=0$, and let  $\psi\colon [0,\infty)\to [0,\infty)$ be a function. 
    \begin{itemize}
        \item $\psi$ is called an \emph{admissible profile} if it is convex, nondecreasing, locally absolutely continuous, satisfies $\psi(0)=0$, and 
        \begin{equation}\label{Admissible}
        M_{\psi} := \int_{0}^{\infty} r^{n-1}e^{-\psi(r)/\delta}\,dr <\infty.
        \end{equation}
        
        \item Given an admissible profile $\psi$, the \emph{comparison radial measure} is the probability measure  $\nu_{\psi,\delta}$ on $[0,\infty)$ whose density with respect to
    Lebesgue measure is
        $$
        \frac{d\nu_{\psi,\delta}}{dr}(r)=\frac{r^{n-1}e^{-\psi(r)/\delta}}{M_{\psi}}.
        $$
        We fix once and for all a random variable $\widetilde R$ with law $\nu_{\psi,\delta}$; it implies that  $\mathbb{E}[\widetilde R]=\int_0^\infty r\,d\nu_{\psi,\delta}(r)$ is the one-dimensional comparison constant that appears in the bounds below.
        \item   We say that $G$ satisfies the \emph{radial $\psi$-dominance condition} if
        \begin{equation}\label{dominance}
        G(r_2\omega)-G(r_1 \omega)\geq \psi(r_2)-\psi(r_1) \quad \textrm{ for all } 0\leq r_1\leq r_2 \textrm{ and } \omega \in \mathbb{S}^{n-1}.
        \end{equation}
    \end{itemize}
\end{definition}

\begin{theorem}\label{Stochastic-dominance}
    Let $G\colon \mathbb{R}^n \to \mathbb{R}\cup \{+\infty\}$ be a proper lower semicontinuous function with unique minimizer at $0$, and assume that $G(0)=0$. Let $\psi$ be an admissible profile satisfying \eqref{Admissible}, and suppose $G$ satisfies the radial $\psi$-dominance condition\eqref{dominance}. For $\delta>0$, define the probability measure \(\mu_\delta\) on
\(\mathbb{R}^n\) by
$$
d\mu_{\delta}(z):=\frac{e^{-G(z)/\delta}}{\int_{\mathbb{R}^n}e^{-G(z)/\delta}\,dz}\,dz, \quad z\in \mathbb{R}^n.
$$
Then the following assertions hold:
\begin{enumerate}
    \item[(i)] The measure $\mu_{\delta}$ is well-defined, and $\mathbb{E}_{\mu_{\delta}}[\Vert Z\Vert]<\infty$.
    \item[(ii)] The mean $\mathfrak{m}_\delta(G):=\mathbb{E}_{\mu_\delta}[Z]$ satisfies
    \begin{equation*}
        \|\mathfrak{m}_\delta(G)\|
        \le
        \mathbb{E}[\widetilde R]
        =
        \frac{\displaystyle \int_0^\infty r^n e^{-\psi(r)/\delta}\,dr}
        {\displaystyle \int_0^\infty r^{n-1} e^{-\psi(r)/\delta}\,dr}.
    \end{equation*}
    \item[(iii)] For every \(\varepsilon>0\), there exists a proper lower
    semicontinuous function \(G_\varepsilon\) satisfying
    \eqref{dominance}  such that
    \[
    \|\mathfrak{m}_\delta(G_\varepsilon)\|
    >
    \mathbb{E}[\widetilde R]-\varepsilon.
    \]
\end{enumerate}
\end{theorem}

\begin{lemma}\label{polar-decomposition}
Under the assumptions of Theorem \ref{Stochastic-dominance}, let \(R:=\|Z\|\) with \(Z\sim\mu_\delta\). Then \(R\) has density
\[
p_R(r)=\frac{r^{n-1}A(r)}{\displaystyle\int_0^\infty s^{n-1}A(s)\,ds},
\qquad
A(r):=\int_{\mathbb S^{n-1}} e^{-G(r\omega)/\delta}\,d\omega,
\qquad r\ge 0.
\]
\end{lemma}

\begin{lemma}\label{likelihood-ratio}
Assume that \eqref{dominance} holds. Then the function
\[
\Lambda(r):=\frac{A(r)}{\omega_{n-1}e^{-\psi(r)/\delta}}
=
\frac{1}{\omega_{n-1}}\int_{\mathbb{S}^{n-1}} e^{-H(r\omega)/\delta}\,d\omega,
\qquad r\ge 0,
\]
is nonincreasing.
\end{lemma}

\begin{proposition}\label{stochastic-dominance-lemma}
Under the hypotheses of Theorem \ref{Stochastic-dominance}, let
\(Z\sim\mu_\delta\), set \(R:=\|Z\|\), and let \(\widetilde R\sim\nu_{\psi,\delta}\).
Then \(R\le_{\textrm{st}}\widetilde R\). In particular, \(\mathbb{E}[R]<\infty\) and
\[
\mathbb{E}_{\mu_\delta}[\|Z\|]
=
\mathbb{E}[R]
\le
\mathbb{E}[\widetilde R].
\]
\end{proposition}

\begin{proof}
By Lemma \ref{polar-decomposition}, the density of \(R\) is \(p_R\), while the density of \(\widetilde R\sim\nu_{\psi,\delta}\) is \(p_{\widetilde R}\), where
\[
p_R(r)=\frac{r^{n-1}A(r)}{\displaystyle\int_0^\infty s^{n-1}A(s)\,ds},
\quad
p_{\widetilde R}(r)=
\frac{r^{n-1}\omega_{n-1}e^{-\psi(r)/\delta}}
{\displaystyle\int_0^\infty s^{n-1}\omega_{n-1}e^{-\psi(s)/\delta}\,ds},
\qquad r\ge 0.
\]
Moreover, by Lemma \ref{likelihood-ratio}, the function $r\mapsto \Lambda(r):=\frac{A(r)}{\omega_{n-1}e^{-\psi(r)/\delta}}$ is nonincreasing on \([0,\infty)\).

Fix \(t\ge 0\). For every \(0\le r_1<t<r_2\), the monotonicity of \(\Lambda\) yields $\Lambda(r_1)\ge \Lambda(r_2)$,  that is,
\[
A(r_1)e^{-\psi(r_2)/\delta}\ge A(r_2)e^{-\psi(r_1)/\delta}.
\]
Multiplying by \((r_1r_2)^{n-1}\) and integrating over
\(\{(r_1,r_2):0\le r_1<t<r_2<\infty\}\), we obtain
\[
\int_0^t r^{n-1}e^{-\psi(r)/\delta}\,dr
\int_t^\infty r^{n-1}A(r)\,dr
\le
\int_0^t r^{n-1}A(r)\,dr
\int_t^\infty r^{n-1}e^{-\psi(r)/\delta}\,dr.
\]
Dividing by
\[
\left(\int_0^\infty r^{n-1}A(r)\,dr\right)
\left(\int_0^\infty r^{n-1}e^{-\psi(r)/\delta}\,dr\right),
\]
it follows that
\[
\mathbb{P}(\tilde{R} \leq t) \mathbb{P}(R>t)\le \mathbb{P}(\widetilde R>t) \mathbb{P}(R\leq t),
\]
which is equivalent to
\[
\mathbb{P}(R>t)\le \mathbb{P}(\widetilde R>t).
\]
Hence \(R\le_{\textrm{st}}\widetilde R\).  Moreover, since \(R\) and \(\widetilde R\) are nonnegative and integrable,
\[
\mathbb{E}[R]
=
\int_0^\infty \mathbb{P}(R>t)\,dt
\le
\int_0^\infty \mathbb{P}(\widetilde R>t)\,dt
=
\mathbb{E}[\widetilde R].
\]
Finally, \(\mathbb{E}[R]=\mathbb{E}_{\mu_\delta}[\|Z\|]\), which gives the conclusion.
\end{proof}

\end{document}
